\section{Introduction}\label{sec:introduction}
Let $n\ge4$ and $1\le d\le n-2$ with $nd$ even. Write $\Omega=\Omega_{n,d}$ for
the set of simple $d$-regular graphs on the vertex set $[n]=\{1,\ldots,n\}$, put
$M=|\Omega|$, and let $\pi$ be the uniform probability measure on $\Omega$. A
\emph{Curveball trade} at a vertex pair $\{x,w\}$ redistributes the vertices
adjacent to exactly one of $x$ and $w$ between these two endpoints, keeping
both endpoint degrees and every other adjacency~\cite{SNBFS14,CBS18}.

We use the following precise discrete-time convention, which is that
of~\cite{S26}. The Curveball chain $K$ chooses an unordered vertex pair
uniformly and moves to a graph chosen uniformly among those obtainable from
the current one by a Curveball trade at that pair, the current graph included.
No additional holding probability is imposed, and time counts trades. For an
unordered pair $a$, let $E_a$ average a function over the graphs obtainable by
a trade at $a$, and put
\begin{equation}\label{eq:operators}
 h_a=I-E_a,\qquad H=\sum_{a\in\binom{[n]}2}h_a,\qquad
 K=I-\frac{H}{\binom n2}.
\end{equation}
The kernel $K$ is symmetric, and $\operatorname{gap}(K)$ is one minus its
second largest eigenvalue; $\operatorname{gap}(H)$ is the least eigenvalue of
$H$ on the orthogonal complement $\mathbf1^\perp$ of the constants in
$L^2(\pi)$. Thus $\operatorname{gap}(K)=\operatorname{gap}(H)/\binom n2$. The
symmetric group $S_n$ acts on $\Omega$ by relabeling vertices, and on functions
by $(T_\sigma f)(G)=f(\sigma^{-1}G)$. Let
\[
 L_T=\sum_{a\in\binom{[n]}2}(I-T_a)
\]
be the Laplacian of the random-transposition walk on $\Omega$, where $T_a$ is
the relabeling by the transposition of the two vertices of $a$. For a
partition $\lambda$ of $n$, let $U_\lambda$ be the $\lambda$-isotypic component
of $L^2(\pi)$ and let
$\operatorname{cont}(\lambda)=\sum_i\bigl[\binom{\lambda_i}2-(i-1)\lambda_i\bigr]$
be the sum of the contents of its cells. A function on $\Omega$ is a
\emph{rooted-class function} if it is invariant under every relabeling that
fixes the vertex $n$; such a function depends only on the isomorphism class of
the graph rooted at $n$. Here and throughout, logarithms are natural.

\begin{theorem}[Regular graphs]\label{thm:main}
Let $n\ge4$ and $1\le d\le n-2$ with $nd$ even.
\begin{enumerate}
\item For every pair $\{x,y\}$, the centered edge indicator
$g_{xy}=\mathbf1[xy]-d/(n-1)$ is a nonzero function with $Hg_{xy}=(n-1)g_{xy}$.
Consequently
\begin{equation}\label{eq:upper}
 \operatorname{gap}(H)\le n-1,\qquad \operatorname{gap}(K)\le\frac2n .
\end{equation}
\item $H\succeq\tfrac12L_T$. Consequently, for every partition $\lambda$ of
$n$,
\begin{equation}\label{eq:sector}
 H\succeq\frac12\Bigl[\binom n2-\operatorname{cont}(\lambda)\Bigr]
 \quad\text{on }U_\lambda .
\end{equation}
The right side is at least $n-1$ for every $\lambda$ other than $(n)$ and
$(n-1,1)$, and equals $n/2$ for $\lambda=(n-1,1)$.
\item Let $\mu$ be the least eigenvalue of $H$ on the centered rooted-class
functions, with $\mu=\infty$ when every rooted-class function is constant.
Then
\begin{equation}\label{eq:reduction}
 \operatorname{gap}(H)=\min\{n-1,\mu\}.
\end{equation}
\item If $d\in\{1,n-2\}$, then $\operatorname{gap}(H)=n-1$ and
$\operatorname{gap}(K)=2/n$.
\end{enumerate}
\end{theorem}

Part~(2) is stated for regular degree sequences. For other sequences the
transposition of two vertices of different degrees does not preserve
$\Omega_d$, and Remark~\ref{cmp:general} gives the comparison for the pairs of
equal degree. No lower bound for $\mu$ beyond the universal bound
$H\succeq2(I-\Pi)$ of~\cite[Theorem~1.1 and Corollary~6.4]{S26} is proved
here for general $n$ and $d$. For $d=1$ the chain is a lazy version of the
walk of Diaconis and Holmes on perfect matchings,
$K=(1-\beta)I+\beta K_{\mathrm{DH}}$ with $\beta=(n-2)/(2(n-1))$, so
part~(4) for $d=1$ follows from their
spectrum~\cite[Theorem~2 and Corollary~1]{DH02}
(Remark~\ref{root:matchings}); for $d=n-2$ it follows from the case $d=1$ by
complementation.

\begin{theorem}[Small cases]\label{thm:small}
Let $1\le d\le n-2$ with $nd$ even. If $4\le n\le\NMAX$, or if $n=12$ and
$d\notin\{5,6\}$, then $\operatorname{gap}(H)=n-1$ and
$\operatorname{gap}(K)=2/n$ on $\Omega_{n,d}$. The lower bound
uses the exact integer computation of Section~\ref{cert:section}.
\end{theorem}

\begin{corollary}[Relaxation and mixing]\label{cor:mixing}
Let $d(s)=\max_{G\in\Omega}\frac12\sum_{G'\in\Omega}|K^s(G,G')-\pi(G')|$ and
$t_{\mathrm{mix}}(\varepsilon)=\min\{s\ge0:d(s)\le\varepsilon\}$. Under the
hypotheses of Theorem~\ref{thm:main}, the relaxation time
$1/\operatorname{gap}(K)$ is at least $n/2$, and for $0<\varepsilon<1/2$,
\begin{equation}\label{eq:mixlower}
 t_{\mathrm{mix}}(\varepsilon)\ge\Bigl(\frac n2-1\Bigr)\log\frac1{2\varepsilon}.
\end{equation}
Whenever $\operatorname{gap}(K)=2/n$, in particular under the hypotheses of
Theorem~\ref{thm:main}(4) or of Theorem~\ref{thm:small}, also
\begin{equation}\label{eq:mixupper}
 t_{\mathrm{mix}}(\varepsilon)\le\Bigl\lceil\frac n2
 \Bigl(\log\frac1{2\varepsilon}+\frac12\log(M-1)\Bigr)\Bigr\rceil .
\end{equation}
\end{corollary}

No sharpness of \eqref{eq:mixlower} or \eqref{eq:mixupper} is asserted.

\paragraph{History and significance.}
Strona, Nappo, Boccacci, Fattorini and San-Miguel-Ayanz~\cite{SNBFS14}
introduced Curveball trades for binary matrices with fixed margins;
Verhelst~\cite{Verhelst08} had developed trade-based sampling for the same
state space. Carstens, Berger and Strona defined undirected
Curveball trades and proved convergence to the uniform
distribution~\cite[Theorem~11]{CBS18}; they left theoretical mixing bounds
open. Carstens and Kleer~\cite{CK18} bounded the relaxation time of the
Curveball chain for binary matrices relative to that of the switch chain;
their comparison gives absolute bounds only where switch-chain bounds are
known, and it does not treat undirected graphs.
Caputo~\cite{Caputo08} studied binary collision processes with a conservation
law, of which row trades on binary matrices with fixed column sums are an
instance, by expanding the square of the generator over triples of particles.
Fu, Qin and Wang~\cite[Lemma~4.1 and Theorem~5.1]{FQW26} used this expansion
to prove that the row-trade chain on binary matrices with arbitrary margins
has gap at least $1/\binom m2$, where $m$ is the number of rows. The
manuscript~\cite{OAI26}, which has not been peer reviewed and whose main
mixing theorem has an agent-written Lean formalization upstream, developed
the pair-fiber operators for undirected graphs.
For the switch chain on regular graphs, Cooper, Dyer and
Greenhill~\cite{CDG07,CDG12} proved polynomial mixing, and Tikhomirov and
Youssef~\cite{TY23} proved sharp Poincar\'e bounds for regular bipartite graphs
in a specified degree range. For every graphical degree sequence with $M>1$,
\cite[Theorem~1.1]{S26} gives $\operatorname{gap}(K)\ge2/\binom n2$, with
equality whenever $M=2$; how much larger the gap is for other
sequences was left open there. Theorem~\ref{thm:main} shows that on regular
graphs the gap is at most $2/n$, that $n-1$ is an eigenvalue of $H$ for every
$n$ and $d$, and that every eigenvalue of $H$ below $n-1$ on $\mathbf1^\perp$,
if one exists, has a rooted-class eigenfunction. Theorem~\ref{thm:small}
excludes such eigenvalues for $n\le\NMAX$ and for most degrees at $n=12$. Whether
$\operatorname{gap}(K)=2/n$ for every regular degree sequence is open
(Section~\ref{op:section}).

\paragraph{Proof strategy.}
On a regular fiber the two endpoints of a trade have the same number of
singleton neighbors, so a trade assigns each singleton to either endpoint with
probability $1/2$. Summing this identity over pairs gives part~(1). The same
balance shows that exchanging the labels of the two endpoints is itself a
trade. Hence every function fixed by $E_a$ is fixed by $T_a$, which gives
$h_a\succeq\tfrac12(I-T_a)$ and part~(2); the class sum of transpositions acts
on $U_\lambda$ as the scalar $\operatorname{cont}(\lambda)$, and an elementary
bound on contents gives the constants. The two components not covered by
part~(2) are generated by the functions fixed by a point stabilizer, which
gives part~(3). For Theorem~\ref{thm:small} we compute the action of $H$ on
rooted-class indicators in rational arithmetic and decide the resulting matrix
inequality by an integer congruence certificate. The comparison with the random transpositions of
vertex labels is the new step; the pair-fiber operators and the kernel of $H$
are taken from~\cite{OAI26,S26}.

\paragraph{Organization.}
Section~\ref{sec:setup} fixes the operators and the relabeling action.
Section~\ref{edg:section} proves the eigenvalue identity.
Section~\ref{cmp:section} proves the comparison with random transpositions and
the bounds on isotypic components. Section~\ref{root:section} proves the
reduction to rooted-class functions, Theorem~\ref{thm:main} and
Corollary~\ref{cor:mixing}. Section~\ref{cert:section} describes the exact
computation behind Theorem~\ref{thm:small}. Section~\ref{exp:section} locates
the slack in the expansion of $H^2$ used in~\cite{S26} and shows that it cannot
be removed term by term. Section~\ref{bin:section} records the corresponding
statements for binary matrices with equal row sums, where the equality
corresponding to $\operatorname{gap}(K)=2/n$ fails for some margins with
equal column sums, and
Section~\ref{op:section} states open problems. All proof ingredients are
supplied in the paper, except the exact computations of
Section~\ref{cert:section}, of Proposition~\ref{exp:obstruction}(2) and of
the binary-matrix examples of Section~\ref{bin:section}, which are carried
out by ancillary programs.
